\section*{Problem 1}

设 $g : \mathbb{R}^2 \times \mathbb{R}^2 \to \mathbb{R}$ 定义为
\[
g \left(\boldsymbol{x}, \boldsymbol{y}\right) = 2 x_1 y_1 + 3 x_1 y_2 - x_2 y_1 + 4 x_2 y_2,
\]
其中 $\boldsymbol{x} = \left(x_1, x_2\right), \boldsymbol{y} = \left(y_1, y_2\right)$.
\begin{enumerate}
    \item 写出 $g$ 在标准基下的 Gram 矩阵 $G$；
    \item 判断 $g$ 是否为对称双线性型.
\end{enumerate}

\begin{proof}
为求解标准基 $\left\{\boldsymbol{e}_1, \boldsymbol{e}_2\right\}$ 下的 Gram 矩阵，代入并求值有
\[
G = \begin{bmatrix}
    g \left(\boldsymbol{e}_1, \boldsymbol{e}_1\right) & g \left(\boldsymbol{e}_1, \boldsymbol{e}_2\right) \\
    g \left(\boldsymbol{e}_2, \boldsymbol{e}_1\right) & g \left(\boldsymbol{e}_2, \boldsymbol{e}_2\right) \\
\end{bmatrix},
\begin{cases}
    g \left(\boldsymbol{e}_1, \boldsymbol{e}_1\right) &= 2, \\
    g \left(\boldsymbol{e}_1, \boldsymbol{e}_2\right) &= 3, \\
    g \left(\boldsymbol{e}_2, \boldsymbol{e}_1\right) &= -1, \\
    g \left(\boldsymbol{e}_2, \boldsymbol{e}_2\right) &= 4. \\
\end{cases}
\implies G = \boxed{\begin{bmatrix}
    2 & 3 \\
    -1 & 4 \\
\end{bmatrix}}.
\]
由此亦可见 $\boxed{g \textrm{ 不是对称双线性型}}$.

\end{proof}

\section*{Problem 2}

设 $V$ 是 2 维线性空间，基为 $\left\{\boldsymbol{e}_1, \boldsymbol{e}_2\right\}$，定义双线性型 $g$ 满足
\[
\begin{aligned}
g \left(\boldsymbol{e}_1, \boldsymbol{e}_1\right) = 1, \\
g \left(\boldsymbol{e}_1, \boldsymbol{e}_2\right) = 2, \\
g \left(\boldsymbol{e}_2, \boldsymbol{e}_1\right) = 3, \\
g \left(\boldsymbol{e}_2, \boldsymbol{e}_2\right) = 4. \\
\end{aligned}
\]
\begin{enumerate}
    \item 写出 $g$ 在基 $\left\{\boldsymbol{e}_1, \boldsymbol{e}_2\right\}$ 下的 Gram 矩阵；
    \item 若新基 $\left\{\boldsymbol{f}_1, \boldsymbol{f}_2\right\}$ 满足 $\boldsymbol{f}_1 = \boldsymbol{e}_1 + \boldsymbol{e}_2, \boldsymbol{f}_2 = \boldsymbol{e}_1 - \boldsymbol{e}_2$，求 $g$ 在新基下的 Gram 矩阵.
\end{enumerate}

\begin{proof}
易见
\[
G = \begin{bmatrix}
    g \left(\boldsymbol{e}_1, \boldsymbol{e}_1\right) & g \left(\boldsymbol{e}_1, \boldsymbol{e}_2\right) \\
    g \left(\boldsymbol{e}_2, \boldsymbol{e}_1\right) & g \left(\boldsymbol{e}_2, \boldsymbol{e}_2\right) \\
\end{bmatrix} = \boxed{\begin{bmatrix}
    1 & 2 \\
    3 & 4 \\
\end{bmatrix}}.
\]
欲求 $g$ 在 $\left\{\boldsymbol{f}_1, \boldsymbol{f}_2\right\}$ 下的 Gram 矩阵，求值
\[
\begin{cases}
    g \left(\boldsymbol{f}_1, \boldsymbol{f}_1\right) = g \left(\boldsymbol{e}_1 + \boldsymbol{e}_2, \boldsymbol{e}_1 + \boldsymbol{e}_2\right) = 10, \\
    g \left(\boldsymbol{f}_1, \boldsymbol{f}_2\right) = g \left(\boldsymbol{e}_1 + \boldsymbol{e}_2, \boldsymbol{e}_1 - \boldsymbol{e}_2\right) = -2, \\
    g \left(\boldsymbol{f}_2, \boldsymbol{f}_1\right) = g \left(\boldsymbol{e}_1 - \boldsymbol{e}_2, \boldsymbol{e}_1 + \boldsymbol{e}_2\right) = -4, \\
    g \left(\boldsymbol{f}_2, \boldsymbol{f}_2\right) = g \left(\boldsymbol{e}_1 - \boldsymbol{e}_2, \boldsymbol{e}_1 - \boldsymbol{e}_2\right) = 0. \\
\end{cases}
\]
于是有
\[
G^\prime = \begin{bmatrix}
    g \left(\boldsymbol{f}_1, \boldsymbol{f}_1\right) & g \left(\boldsymbol{f}_1, \boldsymbol{f}_2\right) \\
    g \left(\boldsymbol{f}_2, \boldsymbol{f}_1\right) & g \left(\boldsymbol{f}_2, \boldsymbol{f}_2\right) \\
\end{bmatrix} \\
= \boxed{\begin{bmatrix}
    10 & -2 \\
    -4 & 0 \\
\end{bmatrix}}.
\]

\end{proof}
\section*{Problem 3}

设 $g \left(\boldsymbol{x}, \boldsymbol{y}\right) = \boldsymbol{x}^\top A \boldsymbol{y}$ 是 $\mathbb{R}^n$ 上的双线性型，$A$ 为 $n \times n$ 矩阵. 证明：$g$ 是对称双线性型，当且仅当 $A$ 是对称矩阵.

\begin{proof}
先证必要性。若 $g$ 为对称双线性型，则
\[
\begin{aligned}
    \boldsymbol{x}^\top A \boldsymbol{y} &= g \left(\boldsymbol{x}, \boldsymbol{y}\right) \\
    &= g \left(\boldsymbol{y}, \boldsymbol{x}\right) = g \left(\boldsymbol{y}, \boldsymbol{x}\right) ^\top = \left(\boldsymbol{y}^\top A \boldsymbol{x}\right)^\top = \boldsymbol{x}^\top A^\top \boldsymbol{y}.
\end{aligned}
\]
令 $\boldsymbol{x}, \boldsymbol{y}$ 取遍 $\left\{\boldsymbol{e}_1, \boldsymbol{e}_2, \cdots, \boldsymbol{e}_n\right\}$，从而
\[
A_{i, j} = \boldsymbol{e}_i^\top A \boldsymbol{e}_j = \boldsymbol{e}_i^\top A^\top \boldsymbol{e}_j = A^\top_{i, j},
\]
故 $A=A^\top$。

再证充分性。若 $A=A^\top$，则对任意 $\boldsymbol{x},\boldsymbol{y}\in\mathbb{R}^n$，
\[
g(\boldsymbol{y},\boldsymbol{x})
=\boldsymbol{y}^\top A\boldsymbol{x}
=\left(\boldsymbol{y}^\top A\boldsymbol{x}\right)^\top
=\boldsymbol{x}^\top A^\top\boldsymbol{y}
=\boldsymbol{x}^\top A\boldsymbol{y}
=g(\boldsymbol{x},\boldsymbol{y}).
\]
因此 $g$ 是对称双线性型。

\end{proof}
